\section{Hopf monoids}\label{s:hopf}


We review basic notions on Hopf monoids and illustrate definitions with three classical examples. We encourage the reader to see~\cite{AM:2010} for a deeper study on this topic. Throughout the paper $\field$ denotes an arbitrary field and all vector spaces are assumed to be over $\field$. 
In general, a Hopf monoid is defined in a symmetric monoidal category, but here we will use the term {Hopf monoid} in a much more
restrictive manner. From now on, \emph{Hopf monoid} stands for a connected Hopf monoid in the symmetric monoidal category of vector species using Cauchy product. Rather than defining these concepts in their full generality, we give here a brief description of the data that is needed for our purposes. However we highly encourage the reader to see~\cite{AM:2010} for more details on Hopf monoids and~\cite{species98} for more details on species.


%%%%%%%%%%%%%%
%% Basic Definition of Hopf monoid
%%%%%%%%%%%%%%
\subsection{Species and Hopf monoids}\label{ss:sp-hopf}
A \emph{vector species} is a functor from the category of finite sets and bijections to the category of vector spaces and linear maps. Informally, a {vector species} $\bH$ is a collection of vector spaces $\bH[I]$,
one for each finite set $I$, equivariant with respect to bijections $I\cong J$.
A morphism of species $f:\bH\to\mathbf{Q}$ is a collection of linear maps
$f_I:\bH[I]\to{\mathbf{Q}}[I]$ which commute with bijections.

A \emph{set composition} of a finite set $I$ is a finite sequence $(A_1,\ldots,A_k)$
of disjoint non-empty subsets of $I$ whose union is $I$. In this situation, we write
$
(A_1,\ldots,A_k)\models I.
$

A \emph{Hopf monoid} consists of a vector species $\bH$
equipped with two collections $\mu$ and $\Delta$ of equivariant linear maps
\[
\bH[A_1]\otimes\bH[A_2] \map{\mu_{A_1,A_2}}\bH[I]
\qand
\bH[I]\map{\Delta_{A_1,A_2}} \bH[A_1]\otimes\bH[A_2]
\]
subject to a number of axioms, of which the main ones follow.

%\noindent
%{\bf Associativity}:
\begin{associativity}
For each set composition $(A_1, A_2, A_3)\models I$, the diagrams
\begin{equation}\label{e:assoc}
\begin{gathered}
\xymatrix@R+1pc@C+20pt{
\bH[A_1]\otimes\bH[A_2]\otimes\bH[A_3]\ar[r]^-{\id\otimes\mu_{A_1,A_2}}
\ar[d]_{\mu_{A_1,A_2}\otimes\id} & \bH[A_1]\otimes\bH[A_2\cup A_3]\ar[d]^{\mu_{A_1,A_2\cup A_3}} \\
\bH[A_1\cup A_2]\otimes \bH[A_3]\ar[r]_-{\mu_{A_1\cup A_2,A_3}} & \bH[I]
}
\end{gathered}
\end{equation}

\begin{equation}
\label{e:coassoc}
\begin{gathered}
\xymatrix@R+1pc@C+40pt{
\bH[I]\ar[r]^-{\Delta_{A_1\cup A_2,A_3}} \ar[d]_{\Delta_{A_1,A_2\cup A_3}}
& \bH[A_1\cup A_2]\otimes \bH[A_3]\ar[d]^{\Delta_{A_1,A_2}\otimes\id}
\\
\bH[A_1]\otimes\bH[A_2\cup A_3]\ar[r]_-{\id\otimes\Delta_{A_1,A_2}}
& \bH[A_1]\otimes\bH[A_2]\otimes\bH[A_3]
}
\end{gathered}
\end{equation}
commute.
\end{associativity}

%\noindent
%{\bf Compatibility}:

\begin{compatibility}
Fix two set compositions $(A_1,A_2)$ and $(B_1,B_2)$ of $I$,
and consider the resulting pairwise intersections:
\[
P:=A_1\cap B_1,\ Q:=A_1\cap B_2,\ R:=A_2\cap B_1,\ T:=A_2\cap B_2,
\]
as illustrated below
\begin{equation}\label{e:4sets}
\begin{gathered}
\quad \begin{picture}(100,90)(20,0)
 \put(50,40){\oval(90,70)}
 \put(5,40){\dashbox{2}(90,0){}}
 \put(45,55){$A_1$}
 \put(45,15){$A_2$}
\end{picture}
\quad
\begin{picture}(100,90)(10,0)
 \put(50,40){\oval(90,70)}
 \put(50,5){\dashbox{2}(0,70){}}
 \put(20,35){$B_1$}
 \put(70,35){$B_2$}
\end{picture}
\quad
\begin{picture}(100,90)(0,0)
\put(50,40){\oval(90,70)}
 \put(5,40){\dashbox{2}(90,0){}}
 \put(50,5){\dashbox{2}(0,70){}}
 \put(20,55){$P$}
 \put(70,55){$Q$}
 \put(20,15){$R$}
 \put(70,15){$T$}
\end{picture}
\end{gathered}.
\end{equation}

For any such pair of set compositions, the diagram
\begin{equation}\label{e:comp}
\begin{gathered}
\xymatrix@R+2pc@C-5pt{
\bH[P] \otimes \bH[Q] \otimes \bH[R] \otimes \bH[T] \ar[rr]^{\cong} & &
\bH[P] \otimes \bH[R] \otimes \bH[Q] \otimes \bH[T] \ar[d]^{\mu_{P,R}
\otimes \mu_{Q,T}}\\
\bH[A_1] \otimes \bH[A_2] \ar[r]_-{\mu_{A_1,A_2}}\ar[u]^{\Delta_{P,Q} \otimes
\Delta_{R,T}} & \bH[I] \ar[r]_-{\Delta_{B_1,B_2}} & \bH[B_1] \otimes
\bH[B_2]
}
\end{gathered}
\end{equation}
must commute. The top arrow stands for the map that interchanges the middle factors.

In addition, we require that the Hopf monoid $\bH$ is \emph{connected}, that is,
 $\bH[\emptyset]\cong\field$ and the maps
\[
\xymatrix{
\bH[I]\otimes\bH[\emptyset] \ar@<0.5ex>[r]^-{\mu_{I,\emptyset}} & 
\bH[I] \ar@<0.5ex>[l]^-{\Delta_{I,\emptyset}} 
}
\qqand
\xymatrix{
\bH[\emptyset]\otimes\bH[I] \ar@<0.5ex>[r]^-{\mu_{\emptyset,I}} & 
\bH[I] \ar@<0.5ex>[l]^-{\Delta_{\emptyset,I}} 
}
\]
are the canonical identifications.

The collection $\mu$ is the \emph{product} and the collection $\Delta$
is the \emph{coproduct} of the Hopf monoid $\bH$. For any Hopf monoid $\bH$ the antipode map $S\colon \bH\to\bH$ is computed using Takeuchi's formula (see Section 8.4.2 of~\cite{AM:2010}). More precisely, for any finite set $I$ and a set composition $A=(A_1,\ldots,A_k)\models I$, if $k=1$ we let $\mu_{A_1}=\Delta_{A_1}=\Id_I$ the identity map on $\bH[I]$, and if $k>1$ we let
 $$ \mu_{A_1,\ldots,A_k} = \mu_{A_1,I\setminus A_1}(\Id_{A_1}\otimes\mu_{A_2,\ldots,A_k})\quad\text{and}\quad
 \Delta_{A_1,\ldots,A_k} =(\Id_{A_1}\otimes\Delta_{A_2,\ldots,A_k}) \Delta_{A_1,I\setminus A_1}.
 $$
We then have
\begin{equation}\label{eq:takeuchi}
 S_I = \sum_{k=1}^{|I|}\sum_{(A_1,\ldots,A_k)\models I} (-1)^{k} \mu_{A_1,\ldots,A_k} \Delta_{A_1,\ldots,A_k}
 = \sum_{A\models I} (-1)^{\ell(A)} \mu_A \Delta_A\,.
\end{equation}
A Hopf monoid is (co)commutative if the left (right) diagram below commutes
for all set compositions $(A_1,A_2)\models I$.
\begin{equation}\label{e:comm}
\begin{gathered}
\xymatrix@C-15pt{
\bH[A_1]\otimes\bH[A_2] \ar[rr]^-{\tau_{A_1,A_2}} \ar[rd]_{\mu_{A_1,A_2}} & & \bH[A_2]\otimes\bH[A_1] \ar[ld]^{\mu_{A_2,A_1}} \\
& \bH[I] &
}
\\
\xymatrix@C-15pt{
\bH[A_1]\otimes\bH[A_2] \ar[rr]^-{\tau_{A_1,A_2}} & &\bH[A_2]\otimes\bH[A_1] \\
& \bH[I] \ar[lu]^{\Delta_{A_1,A_2}} \ar[ru]_{\Delta_{A_2,A_1}} &
}
\end{gathered}
\end{equation}
The arrow $\tau_{A_1,A_2}$ stands for
the map that interchanges the factors.

A morphism of Hopf monoids $f:\bH\to\mathbf{Q}$ is a morphism of species that
commutes with $\mu$ and $\Delta$.
\end{compatibility}


%%%%%%%%%%%%%%
%% HOPF MONOID of LINEAR ORDERs
%%%%%%%%%%%%%%
\subsection{The Hopf monoid of linear orders \texorpdfstring{$\bL$ (\cite{AM:2010})}{L ([4])}}\label{ss:order}

For any finite set $I$ let
$\rL[I]$ be the set of all linear orders on $I$. For instance, if $I=\{a,b,c\}$,
\[
\rL[I]=\{abc,\,bac,\,acb,\,bca,\,cab,\,cba\}.
\]
The vector species $\bL$ is such that $\bL[I]:=\field\rL[I]$ is the vector space with basis $\rL[I]$. 

Given $(A_1,A_2)\models I$
and linear orders $\alpha_1,\alpha_2$ on $A_1,A_2$, respectively,
their concatenation $\alpha_1\cdot \alpha_2$ 
is the linear order on $I$ given by $\alpha_1$ followed by $\alpha_2$.
Given a linear order $\alpha$ on $I$ and $P\subseteq I$, the restriction 
 $\alpha|_P$ is the ordering in $P$ given by the order $\alpha$.
These operations give rise to maps
\begin{equation}\label{e:L}
\begin{aligned}
\rL[A_1]\times\rL[A_2] & \to \rL[I] &\qquad\qquad \rL[I] & \to \rL[A_1]\times\rL[A_2]\\
(\alpha_1,\alpha_2) & \to \alpha_1\cdot\alpha_2
& \alpha & \to (\alpha |_{A_1},\alpha |_{A_2}).
\end{aligned}
\end{equation}
Extending by linearity, we obtain linear maps
\[
 \mu_{A_1,A_2}:\bL[A_1]\otimes\bL[A_2] \to \bL[I]
 \qand
\Delta_{A_1,A_2}:\bL[I]\to \bL[A_1]\otimes\bL[A_2]
\]
which turn $\bL$ into a cocommutative but not commutative Hopf monoid. The reader should check that all the required axioms for a Hopf monoid are indeed satisfied for this and the upcoming examples.


%%%%%%%%%%%%%%
%% HOPF MONOID of SET PARTITIONs
%%%%%%%%%%%%%%
\subsection{The Hopf monoid of set partitions \texorpdfstring{$\bPi$ (\cite{AM:2010})}{pi ([4])} }\label{ss:partition}

A \emph{partition} of a finite set $I$ is a collection $X$ of disjoint nonempty subsets whose union is $I$. The subsets are the \emph{blocks} of $X$.

Given a partition $X$ of $I$ and $P\subseteq I$, the restriction
$X|_P$ is the partition of $P$ whose blocks are 
the nonempty intersections of the blocks of $X$ with $P$.
Given $(A_1,A_2)\models I$
and partitions $X_i$ of $A_i$, $i=1,2$, 
the union $X_1\cup X_2$ is the partition of $I$ 
whose blocks are the blocks of $X_1$ and the blocks of $X_2$.

Let ${\bpi}[I]$ denote the set of partitions of $I$ and $\bPi[I]=\field {\bpi}[I]$ the vector
space with basis ${\bpi}[I]$. 
A Hopf monoid structure on $\bPi$ is defined and studied in~\cite{AguiarBergeronThiem,AM:2010,BakerJarvisBergeronThiem,BergeronZabrocki}. Among its various linear bases, we are interested in the \emph{power-sum} basis 
on which the operations are as follows. 
The product
\[
\mu_{A_1,A_2}: \bPi[A_1] \otimes \bPi[A_2] \to \bPi[I] 
\]
is given by
\begin{equation}\label{e:prod-m}
\mu_{A_1,A_2}({X_1} \otimes {X_2}) = X_1\cup X_2.
\end{equation}
for $X_i\in {\bpi}[A_i]$ and extended linearly. The coproduct
\[
\Delta_{A_1,A_2}: \bPi[I] \to \bPi[A_1]\otimes\bPi[A_2]
\]
is given by
\begin{equation}\label{e:coprod-m}
\Delta_{A_1,A_2}(X) = 
\begin{cases}
{X|_{A_1}} \otimes {X|_{A_2}} & \text{if $A_1$ is the union of some blocks of $X$,} \\
0 & \text{otherwise,}
\end{cases}
\end{equation}
for $X\in {\bpi}[I]$ and extended linearly.
These operations turn the species $\bPi$ into a Hopf monoid that is both commutative and cocommutative.

%%%%%%%%%%%%%% APR-
%% HOPF MONOID of GRAPHS
%%%%%%%%%%%%%%
\subsection{The Hopf monoid of simple graphs \texorpdfstring{$\mathbf{G}$}{G}}\label{ss:graph}

A \emph{(simple) graph} $g$ on a finite set $I$ is a collection $E$ of subsets of $I$ of size 2.
% relation on $I$ that is symmetric and irreflexive. 
The elements of $I$ are the \emph{vertices} of $g$.
There is an \emph{edge} between two vertices $i, j$ if $e=\{i,j\}\in E$. In this case we say that $e$ is \emph{incident} to $i$ and $j$.


 Given a graph $g$ on $I$ and $P\subseteq I$, the restriction
$g|_P$ is the graph on the vertex set $P$ whose edges are the edges of $g$ incident to elements of $P$ only.
Let $(A_1, A_2)\models I$ and
consider the graphs $g_i$ of $A_i$, for $i=1,2$.
The union $g_1\cup g_2$ is the graph on $I$ 
whose edges are those of $g_1$ and those of $g_2$.

Let $\rG[I]$ denote the set of graphs on $I$ and $\bG[I]=\field \rG[I]$ the vector
space with basis $\rG[I]$. 
A Hopf monoid structure on $\bG$ is defined using the maps
\begin{equation}\label{e:G}
\begin{aligned}
\rG[A_1]\times\rG[A_2] & \to \rG[I] &\qquad\qquad \rG[I] & \to \rG[A_1]\times\rG[A_2]\\
(g_1,g_2) & \to g_1\cup g_2
& g & \to (g |_{A_1},g |_{A_2}).
\end{aligned}
\end{equation}
Extending linearly, we obtain linear maps
\[
 \mu_{A_1,A_2}:\bG[A_1]\otimes\bG[A_2] \to \bG[I]
 \qand
\Delta_{A_1,A_2}:\bG[I]\to \bG[A_1]\otimes\bG[A_2].
\]
These operations turn the species $\bG$ into a Hopf monoid that is both commutative and cocommutative.

%%%%%%%%%%%%%%
%% HOPF MONOID of HYPERGRAPHS
%%%%%%%%%%%%%%
\subsection{The Hopf monoid of simple hypergraphs \texorpdfstring{$\mathbf{HG}$}{HG}}\label{ss:hypergraph}

Let $2^I$ denote the collection of subsets of $I$. Let $\mathbf{HG} [I]=\field {\mathbf{hg}} [I]$ be the space spanned by the basis $\mathbf{hg} [I]$ where
 $$\mathbf{hg}[I]=\big\{h \subseteq 2^I\,: U\in h\text{ implies }|U|\ge 2\big\}.$$
 %\caro{but then, what is the construction on sets of size 1? it is only the empty set} 
An element $h\in\mathbf{hg}[I]$ is a \emph{hypergraph on $I$}.
For $(P,T)\models I$ and $h,k\in \mathbf{hg}[I]$, the multiplication is given by $\mu_{P,T}(h,k)=h\cup k$ and the comultiplication 
is given by $\Delta_{P,T}(h) = h |_P\otimes h |_T$ where $h |_P=\{U\in h : U\cap P = U\}$. Extending these definitions linearly we have that $\mathbf{HG}$ is 
commutative and cocommutative Hopf monoid. 

%\begin{rema}\label{r:selfdual}
%The dual of a species $\bH$ is the collection $\bH^*$ of dual vector spaces:
%$\bH^*[I]=\bH[I]^*$.
%A species $\bH$ is said to be finite-dimensional if each space $\bH[I]$
%is finite-dimensional.
%Dualizing the operations of a finite-dimensional Hopf monoid $\bH$, 
%one obtains a Hopf monoid $\bH^*$. The Hopf monoid $\bH$ is called self-dual
%if $\bH\cong\bH^*$. In general, such isomorphism is not unique.
%
%Over a field of characteristic $0$, a Hopf monoid that is connected, commutative and cocommutative, is always
%self-dual. This is a consequence of the Cartier-Milnor-Moore theorem.
%The isomorphism with the dual is not canonical.
%In particular, the Hopf monoids $\bPi$ and $\bG$ are self-dual.
%\end{rema}

%%%%%%%%%%%%%%
%% HADAMARD PRODUCT
%%%%%%%%%%%%%%
\subsection{The Hadamard product}\label{ss:hadamard}

Given two species $\bH$ and $\mathbf{Q}$, their \emph{Hadamard product} is the species $\bH\times\mathbf{Q}$ defined by
\[
(\bH\times\mathbf{Q})[I] = \bH[I]\otimes\mathbf{Q}[I],
\]
where $\otimes$ is the usual tensor product of vector spaces over $\field$.
If $\bH$ and $\mathbf{Q}$ are Hopf monoids, then so is $\bH\times\mathbf{Q}$, with the following operations. For $(A_1, A_2)\models I$, the product on $\bH\times\mathbf{Q}$ is depicted in the diagram:
\[
\xymatrix@C+12pt{
(\bH\times\mathbf{Q})[A_1]\otimes(\bH\times\mathbf{Q})[A_2] \ar@{=}[d] \ar@{-->}[rr] &&
 (\bH\times\mathbf{Q})[I] \ar@{=}[dd] \\
\bH[A_1]\otimes \mathbf{Q}[A_1]\otimes \bH[A_2]\otimes \mathbf{Q}[A_2]
\ar[d]_{\cong} && \\
\bH[A_1]\otimes \bH[A_2]\otimes \mathbf{Q}[A_1]\otimes \mathbf{Q}[A_2]
\ar[rr]_-{\mu_{A_1,A_2}^{\mathbf{H}}\otimes\mu_{A_1,A_2}^{\mathbf{Q}}} &&
\bH[I]\otimes\mathbf{Q}[I]
}
\]
The coproduct is defined similarly. If $\bH$ and $\mathbf{Q}$ are (co)commutative, then so is $\bH\times\mathbf{Q}$. 

%%%%%%%%%%%%%%
%% linearized
%%%%%%%%%%%%%%
\subsection{Linearized Hopf monoids}\label{ss:stronglin} A \emph{set species} $\bh$ is a collection of sets $\bh[I]$,
one for each finite set $I$, equivariant with respect to bijections $I\cong J$.
We say that $\bh$ is a basis for a Hopf monoid $\bH$ if for every finite set $I$ we have that $\bH[I]=\field \bh[I]$, the vector space with basis $\bh[I]$. We say that the monoid $\bH$ is \emph{linearized} in the basis $\bh$ if the product and coproduct maps have the following properties.
The product
\[
\mu_{A_1,A_2}: \bH[A_1] \otimes \bH[A_2] \to \bH[I] 
\]
is the linearization of a map
\begin{equation}\label{e:prod-lin}
\mu_{A_1,A_2}: \bh[A_1] \times \bh[A_2] \to \bh[I] 
\end{equation}
and the coproduct
\[
\Delta_{A_1,A_2}: \bH[I] \to \bH[A_1]\otimes\bH[A_2]
\]
is the linearization of a map
\begin{equation}\label{e:coprod-lin}
\Delta_{A_1,A_2}: \bh[I] \to (\bh[A_1]\times\bh[A_2])\cup\{0\}\,.
\end{equation}
It is understood that for any bijection $\sigma\colon I \to J$, the linear isomorphism $\bH[\sigma]\colon \bH[I]\to \bH[j]$ maps the basis $\bh[I]$ to the basis $\bh[J]$ via the bijection $\bh[\sigma]$.
From now on, we will use capital letters for vector species and lower case for set species.

The Hopf monoids $\bL$, $\bPi$, $\bG$ and $\mathbf{HG}$ are linearized in the bases $\mathbf{l}$, $\bpi$, $\mathbf{g}$ and $\mathbf{hg}$ respectively. As remarked in~\cite{MarbergLin}, many of the Hopf monoids in the literature are linearized in some basis. 
%On the other hand, the Hopf monoid $\bL^{\star}$ is not linearized by $\bf l^{*}$, where $\bL^{\star}$ denotes the Hopf monoid \emph{dual} to $\bL$ (see~\cite{AM:2010}).

%%%%%%%%%%%%%%
%% FUNCTOR K and Kbar
%%%%%%%%%%%%%%
\subsection{Functors \texorpdfstring{$\mathcal{K}$}{K} and \texorpdfstring{$\overline{\mathcal{K}}$}{overline K}} \label{ss:K_Kbar} 
As describe in~\cite[Section~III]{AM:2010}, there are some interesting functors from the category of species to the category of graded vector spaces. Let $[n]:=\{1,2,\ldots,n\}$ and assume throughout that $\text{char}(\field)=0$. Given a species $\bH$, we write $\bH[n]$ instead of $\bH[[n]]$. The symmetric group $S_n$ acts on $\bH[n]$ by relabelling. Define the functors $\Kc$ and $\Kcb$ as
 $$ \Kc(\mathbf{H}) = \bigoplus_{n\ge 0} \mathbf{H}[n]
 \qquad\qquad
 \Kcb(\mathbf{H}) = \bigoplus_{n\ge 0} \mathbf{H}[n]_{S_n}
 $$
where
 $$
 \mathbf{H}[n]_{S_n}={\bH}[n]\Big/ \langle x-{\bH}[\sigma](x) \mid \ \sigma\in S_n; \ x\in {\bH}[n]\rangle
 $$
 denotes the quotient space of equivalence classes under the $S_n$-action. When $\bH$ is a Hopf monoid, 
 we can build a product and coproduct
on $ \Kc({\bH})$ and $ \Kcb({\bH})$ from those of $\bH$ together
with certain canonical transformations.
For example, one has that
\[
\Kcb(\bL)\cong\field[t]
\]
is the polynomial algebra on one generator, while $\Kc(\bL)$
is the Hopf algebra introduced by Patras and Reutenauer in~\cite{PR:2004}. The antipode map $S:\bL\rightarrow\bL$ is such that
for $\alpha=a_1\cdots a_n\in\mathbf{l}[n]$
 $$ S_{[n]}(\alpha)=(-1)^n a_n\cdots a_1.$$
 However, the antipode of the graded Hopf algebra $\Kc[\bL]$ is not given by the formula above (see Section~\ref{s:KLxH}). On the other hand, in the Hopf algebra $\Kcb[\bL]\cong \field[t]$, the antipode is given by $S(t^n)=(-1)^nt^n$ and it is the functorial image of the map above. This is not an accident: the functor $\Kc$ may not preserve the antipode but the functor $\Kcb$ always does.

%One has that $\Kcb(\bPi)$ is the 
%Hopf algebra of \emph{symmetric functions}, while 
%$\Kc(\bPi)$ is the Hopf algebra of \emph{symmetric functions in noncommuting variables}, an object studied in various
%references including~\cite{AM:2006,BHRZ:2006,BZ:2009,RS:2006,BakerJBergeronThiem}.

A very interesting relation between the functors $\Kc$ and $\Kcb$ is given in~\cite[Theorem~15.13]{AM:2010} as follows
\begin{equation}\label{eq:KbarK}
 \Kcb(\bL\times\bH) \cong \Kc(\bH),
\end{equation}
where $\bH$ is an arbitrary Hopf monoid. In this paper we aim to make use of this relation to study the antipode problem for some Hopf algebras.


%%%%%%%%%%%%%%
%% ANTIPODE FOR LxP
%%%%%%%%%%%%%%
\section{Antipode for linearized Hopf Monoid \texorpdfstring{$\bL\times \bH$}{L times H}}\label{s:antiLxH}

In this section we show a multiplicity free and cancellation free formula for the antipode of Hopf monoids of the form $\bL\times \bH$ where $\bH$ is linearized in some basis. Thus, by (\ref{eq:KbarK}), we obtain an antipode formula for $\Kc(\bH)$ as well. However, in $\Kc(\bH)$ this antipode formula may not be cancellation free. 

%%%%%%%%%%%%%%
%% ANTIPODE FROM TAKEUCHI
%%%%%%%%%%%%%%
\subsection{Antipode Formula for \texorpdfstring{$\bL\times \bH$}{L times H}}\label{ss:UsingTakeushi}

Let $\bH$ be a Hopf monoid linearized in the basis $\bh$. We intend to resolve the cancellations in the Takeuchi formula for the antipode of $\bL\times\bH$. For a fixed finite set $I$ let $(\alpha,x)\in (\mathbf{l}\times\bh)[I]$, that is, $\alpha$ is a linear ordering on $I$ and $x$ is an element of $\bh [I]$. From~\eqref{eq:takeuchi} we have
\begin{equation}\label{eq:takeuchi1}
 S_I(\alpha,x) = \sum_{A\models I} (-1)^{\ell(A)} \mu_A \Delta_A(\alpha,x)
 = \sum_{\stackrel{A\models I} %\atop 
 {\Delta_A(x)\ne 0}} (-1)^{\ell(A)} (\alpha_A,x_A),
\end{equation}
 summing over all $A=(A_1,\ldots,A_k)\models I$, where $\alpha_A$ denotes the element in $\mathbf{l} [I]$ given~by
 $$\alpha_A= \alpha|_{A_1}\cdot \alpha|_{A_2}\cdots \alpha|_{A_k}\,.
$$ Also, provided $\Delta_A(x)\neq 0$ we set
 $$x_A=\mu_A\Delta_A(x) \in \bh[I]\,.$$
Each composition $A$ gives rise to single elements $\alpha_A$ and $x_A$ since $\bL$ and $\bH$ are linearized in the basis $\mathbf{l}$ and $\bh$, respectively. 
We can thus rewrite equation~\eqref{eq:takeuchi1} as
\begin{equation}\label{eq:takeuchi2}
 S_I(\alpha,x) = \sum_{(\beta,y)\in (\mathbf{l}\times\bh)[I]}\left(\sum_{\stackrel{A\models I}
 % \atop 
 {(\alpha_A,x_A)=(\beta,y)}} (-1)^{\ell(A)}\right) (\beta,y)\,.
\end{equation}
Let 
\begin{equation}\label{eq:alphabetaxy}
 \mathcal{C}_{\alpha,x}^{\beta,y}=\big\{A\models I : (\alpha_A,x_A)=(\beta,y)\big\}\,.
 \end{equation}

Using the notation above we have the following theorem which provides us a multiplicity-free and cancellation-free formula for the antipode of $\bL\times\bH$. 
